\documentclass[10pt,a4paper]{amsart}
\usepackage[margin=19mm]{geometry}
\usepackage{amsmath,amssymb,mathtools}
\usepackage[scr=rsfs,cal=euler]{mathalfa}
\usepackage{xfrac}
\usepackage[unicode,hidelinks,bookmarks=false]{hyperref}
\newcommand{\ie}{i.e.\ }
\newcommand{\cf}{cf.\ }
\newcommand{\resp}{respectively\ }
\newcommand{\Subsection}{\subsection}
\providecommand{\nobreakdash}{\nobreak\hbox{-}}
\setlength{\emergencystretch}{2em}
\multlinegap0pt
\usepackage{amssymb}
\newtheorem{lemma}{Lemma}
\newtheorem{proposition}{Proposition}
\newcommand{\N}{\mathbb{N}}
\newcommand{\setcon}[2]{\left\{#1\ \left|\ #2\right.\right\}}
\DeclareMathOperator{\tw}{tw}
\title{Separation profiles of free products}
\author{David Hume}
\date{Source: arXiv:2604.24462v1 (CC BY 4.0)}
\begin{document}
\maketitle

\noindent\emph{Source and licence.} Separation profiles of free products. Version arXiv:2604.24462v1 (CC BY 4.0), distributed by arXiv under the Creative Commons Attribution 4.0 International licence, http://creativecommons.org/licenses/by/4.0/, as recorded in the arXiv licence metadata for this exact version and shown on the abstract page https://arxiv.org/abs/2604.24462v1. Full licence notice: \texttt{licenses/arxiv-2604.24462-CC-BY-4.0.txt}.\par\smallskip

\noindent\textit{Editorial scope.} Counted unit: the article's proposition proposition (source0.tex lines 296--300). The article's complete original proof of it is delivered with it (source0.tex lines 301--325, one \texttt{proof} environment of 25 lines). No mathematics is written, repaired or re-derived: the statement and the proof are the article's own lines, in the article's own notation, and no retained line is substituted or re-flowed.  Retained uncounted as the source-local context the proof needs: lines 287--289.  Nothing was omitted from any retained span, so the package carries no omission record.  No retained line contains a citation, so the wrapper carries no bibliography and no bibliography entry.  No retained line contains a figure environment, an image-inclusion command or drawing code, so no image is delivered and the package declares no asset dependency.  Editorial formatting only, in addition: an AMS \texttt{amsart} preamble that restates the article's own declarations and macros used by the retained lines, namely the environments \texttt{lemma}, \texttt{proposition} on separate counters, together with the standard packages those lines need.  The retained environments print in this document as Lemma 1, Proposition 1 in order of appearance, whereas the article prints the same items with the numbering of its own full document.  The complete native source of the article as distributed in the arXiv e-print for this exact version is bundled unchanged as \texttt{sources/199226/source0.tex}.
\begin{lemma}\label{lem:twjoin}
    Let $\Gamma=(V\Gamma,E\Gamma)$ be a finite graph admitting two connected subgraphs $\Gamma^1,\Gamma^2$ such that $\Gamma=\Gamma^1\cup\Gamma^2$ and $\Gamma^1\cap \Gamma^2$ is either empty or a single vertex. Then $\tw(\Gamma)=\max\{\tw(\Gamma^1),\tw(\Gamma^2)\}$.
\end{lemma}

\begin{proposition} Let $X$ be a graph and let $\{X_i\}_{i\in I}$ be a family of subgraphs of $X$. If $X$ is tree-graded with respect to $\{X_i\}_{i\in I}$, then
\[
    \tw_X(r)=\max_i \{\tw_{X_i}(r)\}
\]
\end{proposition}

\begin{proof}
    As each $X_i$ is a subgraph of $X$, the first inequality is automatic. Now fix $r\in \N$, let $\Gamma\leq X$ with $|V\Gamma|\leq r$ and define $\Gamma_i=\Gamma\cap X_i$. We will prove that $\tw(\Gamma)\leq\max\setcon{\tw(\Gamma_i)}{i\in I}$. 
    We may assume that $\Gamma$ is connected, as the treewidth of a graph equals the maximum of the treewidths of its connected components. We may also assume that $|V\Gamma|\geq 2$, as the result clearly holds for a single vertex.

    Set $I_\Gamma= \setcon{i\in I}{|V\Gamma_i|\geq 2}$. Note that $I_\Gamma$ is finite. We will carefully verify that $\Gamma$ is tree-graded with respect to $\{\Gamma_i\}_{i\in I_\Gamma}$. 

    Suppose for a contradiction that some $\Gamma_i$ is not connected. Let $P$ be a minimal length path in $X_i$ connecting two vertices $v,w$ from distinct connected components $\Gamma^0_i$ and $\Gamma^1_i$ of $\Gamma_i$. By minimality $P\cap \Gamma=\{v,w\}$. Let $Q$ be any path from $w$ to $v$ in $\Gamma$. The concatenation $PQ$ is a simple cycle. It is not contained in $X_i$, as there is no path $Q$ connecting $w$ to $v$ in $\Gamma_i=\Gamma\cap X_i$ and it is not contained in any other $X_j$ as $P$ contains an edge from $X_i$ and distinct pieces intersect in at most a single vertex, contradicting the definition of a tree-graded graph.

    Next, we show that $\Gamma=\bigcup_{i\in I_\Gamma} \Gamma_i$. As $X=\bigcup_{i\in I} X_i$, $\Gamma=\bigcup_{i\in I} \Gamma_i$, it suffices to show that whenever $|V\Gamma_i|=1$, there is some $\Gamma_j$ containing $\Gamma_i$ with $|V\Gamma_j|\geq 2$. Suppose $\Gamma_i$ is a single vertex $v$. As $\Gamma$ is connected, and we assumed $|V\Gamma|\geq 2$, $v$ has a neighbour $w$ in $\Gamma$. The edge $vw$ is contained in some $X_j$ and is therefore in $\Gamma_j$. It follows that $\{v\}=\Gamma_i\subset \Gamma_j$ and $|V\Gamma_j|\geq 2$ as required.

    The final two properties are very simple. For any $i,j\in I$, $\Gamma_i\cap\Gamma_j\subseteq X_i\cap X_j$ which is either empty or a single vertex. Any simple loop in $\Gamma$ is contained in some $X_i$, so is also contained in $\Gamma_i$.

    Order $I_\Gamma=\{i_0,\ldots,i_k\}$ as follows. Choose $i_0$. For each $m\geq 1$ we choose $i_m$ so that the subgraphs $\bigcup_{j=0}^{m-1} \Gamma_{i_j}$ and $\Gamma_{i_m}$ intersect. This is always possible as $\Gamma=\bigcup_{i\in I_\Gamma} \Gamma_i$ is connected. We claim that this intersection must be a single vertex. Suppose not, and choose distinct $v,w\in \bigcup_{j=0}^{m-1} \Gamma_{i_j}\cap\Gamma_{i_m}$ at minimal distance in $\Gamma_{i_m}$. Let $P$ be a path from $v$ to $w$ in $\Gamma_{i_m}$ of minimal length. It follows from minimality that $P\cap \bigcup_{j=0}^{m-1} \Gamma_{i_j}=\{v,w\}$. let $Q$ be a path from $w$ to $v$ in $\bigcup_{j=0}^{m-1} \Gamma_{i_j}$ of minimal length. The concatenation $PQ$ is a simple cycle that is not contained in a single $\Gamma_i$, contradicting the definition of a tree-grading.
    
    Now we iteratively apply Lemma \ref{lem:twjoin} to $\Gamma_1=\bigcup_{j=0}^{m-1} \Gamma^{i_j}$ and $\Gamma_2=\Gamma^{i_{m}}$ for each $1\leq m\leq k$.

    It follows that
    \[
        \tw(\Gamma) \leq \max_{i\in I_\Gamma} \{\tw(\Gamma_i)\} \leq \max_{i\in I} \{\tw_{X_i}(r)\}
    \]
    As this holds for all subgraphs $\Gamma\leq X$ with $|V\Gamma|\leq r$ we deduce that
    \[
        \tw_X(r) \leq \max_{i\in I} \{\tw_{X_i}(r)\}  \qedhere
    \]
\end{proof}

\end{document}
